%=======================================================================================
% METHODOLOGY CHAPTER -- 2026-08-13, 12:30
% Restructured from methodology_merged_2026-08-11_2221.tex per the structural review notes of
% 2026-08-13. Content (equations, numbers, model list) is unchanged from that version except
% where a fix below specifically required a correction; this pass only reorganises the
% hierarchy, renames headings to match what each part actually contains, and removes code-level
% identifiers (file names, script/class/column names) in favour of plain descriptions. No
% performance results, interpreted associations, or model-comparison verdicts appear here --
% see the Results chapter. Sentences are left terse/skeletal in places deliberately -- fill in
% final prose by hand. Points still flagged "% TODO" were flagged in the prior version and have
% not yet been resolved.
% 2026-08-14: internal \S-cross-references and their \label{sec:...} anchors removed throughout
% (author's own preference, too many labels to manage in Overleaf) -- kept only for equations,
% tables, figures, the appendix, and the other chapter (Data). Section content/order unchanged.
%=======================================================================================

\chapter{Methodology and Experimental Design}
\label{ch:methodology}

This chapter describes how the raw plot data (Chapter~\ref{ch:data}) were converted into the
results used to answer the project's three research questions. It is organised by method, not by
the order experiments were run: the experimental framework first, then common design decisions and
leakage control, environmental feature-set construction, the top-height prediction models, the two
attribution analyses, evaluation design, and finally assumptions/limitations and reproducibility.

\section{Research questions and experimental framework}

The chapter's three research questions split into two analytical modes. RQ1 is prediction: it
establishes whether these models can predict top height competitively, and under what input and
generalisation conditions. One of RQ1's own model families, the pooled Chapman--Richards curve,
also supplies the fitted curve RQ2's own target is built from; RQ3's target is independent of RQ1,
built entirely from each plot's own height/age history and an external yield-class lookup. RQ2 and
RQ3 are both attribution, though: RQ2 pools every plot's residual against that one shared curve,
RQ3 fits each plot its own asymptote from only that plot's own repeated observations.
Table~\ref{tab:rq-framework} summarises all three; each RQ's own section opens with the specific
questions it tests before describing its models.

\begin{table}[!htbp]
\centering
\scriptsize
\renewcommand{\arraystretch}{1.2}
\begin{tabular}{p{0.8cm}p{1.4cm}p{3.3cm}p{4.2cm}p{2.9cm}}
\toprule
RQ & Unit & Target & Model families compared & Evaluation design \\
\midrule
RQ1 & Plot-survey row & 95th-percentile LiDAR top height & Chapman--Richards curve, average-by-age
lookup, linear/random forest/XGBoost, deep neural network, Chapman--Richards-guided neural network
variants & Spatial-block (primary), pooled spatial $k$-fold, temporal, plot-level \\
RQ2 & Plot & Mean residual from the single pooled Chapman--Richards curve, across a plot's survey
history & Elastic Net, XGBoost, two-stage mixed-effects model & Spatial-block \\
RQ3 & Plot & Difference between a plot's own fitted asymptote and its yield-class-implied
asymptote & Elastic Net, XGBoost, GNNWR & Spatial-block \\
\bottomrule
\end{tabular}
\caption{Research questions mapped to unit of analysis, target, model families, and evaluation
design.}
\label{tab:rq-framework}
\end{table}

\section{Common experimental design and leakage control}

\subsection{Spatial and temporal evaluation}

4survey is the primary cohort for every research question, used for the spatial-generalisation
splits below; 6survey (\S\ref{sec:data-cohorts}) is used for the temporal split and run in
parallel elsewhere as a secondary check wherever sample size permits. RQ1's unit of analysis is
the plot-survey row; RQ2 and RQ3 aggregate to one row per plot (Table~\ref{tab:rq-framework}).
Four split types, each answering a different generalisation question:

\begin{itemize}
\setlength\itemsep{0pt}
\item \textbf{Plot level} (naive baseline, not a generalisation claim). Individual plots are split
  at random (60/20/20, seed 42) with no spatial holdout at all (Figure~\ref{fig:split-maps}, top
  row): a held-out plot's nearest neighbour is typically a training plot metres away, so a model
  can partly succeed by picking up on local conditions it was effectively also trained on, rather
  than genuinely generalising. Kept only as an easy-case reference for the splits below to be
  compared against.
\item \textbf{Spatial block} (primary). Addresses that leakage directly: whole forestry
  compartments, not individual plots, are assigned to train/validation/test, so no compartment
  appears on both sides of a split and a model can no longer lean on a near-identical neighbouring
  plot seen during training. Compartments are chosen by row-count-balanced random assignment
  (60/20/20 train/validation/test, seed 42), with a 60\,m buffer around every validation/test
  plot: any training plot within 60\,m of a held-out plot is excluded from training
  (validation/test plots are never removed; Figure~\ref{fig:split-maps}, bottom row, shows the
  resulting assignment and buffer). The buffer distance is set from typical plot spacing rather
  than from a measured spatial-autocorrelation range; leakage protection at a larger spatial scale
  comes from holding out whole compartments, not from the buffer alone
  \citep{roberts2017}. % TODO: citation -- Roberts et al. 2017; no .bib entry yet.
  % A semivariogram of the shared-curve residual was also fitted (test-partition only), giving a
  % range far larger than this buffer -- confirmed NOT used to set the buffer distance, a
  % deliberately separate diagnostic; see the semivariogram mention under Model evaluation,
  % uncertainty and interpretation.
\item \textbf{Spatial block, pooled $k$-fold} (stronger confirmation of the single spatial-block
  split). Compartments are assigned to 5 row-count-balanced folds; each fold in turn is held out
  as test, the next as validation, the remainder as train, so every compartment is evaluated
  exactly once. Predictions from all 5 folds are pooled into one population-wide evaluation set
  rather than averaged as 5 separate summary statistics.
\item \textbf{Temporal} (secondary, harder test). Whole survey years, not compartments, are
  assigned to train/validation/test: 4survey trains on 2008 and 2012, validates on 2021, tests on
  2023; 6survey trains on 2002/2006/2008/2012, validates on 2021, tests on 2023. This asks whether
  a model trained on earlier growth still predicts a real future survey it never saw. This is a
  different, harder question than spatial generalisation (Figure~\ref{fig:temporal-split}). A
  narrow-gap variant (training years closer to the test year) is also defined as a secondary
  check, with validation deliberately assigned to the earliest available year rather than a
  middle year.
  % TODO: the reason for that choice (a failed smoke test with a middle validation year) is
  % implementation history -- move to Appendix or omit from the main text.
  % TODO: no narrow-gap run was found anywhere in TEMP_results as of 2026-08-14 -- confirm this
  % variant has actually been run before citing "its own results" in the Appendix pointer in
  % Reproducibility and implementation, below; if it hasn't been run, soften that pointer or drop
  % it.
\end{itemize}

% FIGURE PLACEHOLDER (plot-level vs. spatial-block split, both cohorts) -- not finalised, likely
% candidate: plot_level_split map vs. spatial_block_split + buffer map, 4survey and 6survey.
% Style TBD -- match whatever map convention is settled on for the rest of the dissertation
% (e.g. the data chapter's cohort coverage map) rather than deciding it here in isolation.
\begin{figure}[!htbp]
\centering
\fbox{\parbox{0.9\textwidth}{\centering\vspace{4cm}FIGURE PLACEHOLDER --- plot-level vs.
spatial-block split maps\vspace{4cm}}}
\caption{Plot-level (random) versus spatial-block (compartment, buffered) split assignment,
4survey and 6survey.}
\label{fig:split-maps}
\end{figure}

% FIGURE PLACEHOLDER (temporal split timeline) -- not finalised.
\begin{figure}[!htbp]
\centering
\fbox{\parbox{0.9\textwidth}{\centering\vspace{3cm}FIGURE PLACEHOLDER --- temporal split
timeline\vspace{3cm}}}
\caption{Temporal split: survey years assigned to train, validation and test, both cohorts.}
\label{fig:temporal-split}
\end{figure}

\paragraph*{Training, validation and test separation.}
Every model in this project reads the same split assignment for a given (cohort, split type,
seed) combination, so results are always compared on identical train/validation/test membership.
Compartment/block/sub-compartment identifier columns are never included in any model's feature
list, enforced automatically. Shared data partitions and matched architectures reduce confounding
when comparing model families under the same split, for example the DNN and the
Chapman--Richards-guided network. Residual differences can still reflect optimisation behaviour or
stochastic variation rather than the model design alone.

\paragraph*{Preprocessing and leakage prevention.}
Continuous features are standardised to zero
mean and unit variance using parameters fitted exclusively on the training partition to prevent
scaling data leakage. Age is scaled independently using a separate scaler, allowing its
training-split standard deviation to be extracted directly for the physics-informed loss
function's unit conversion. Unordered categorical soil variables
(three raster classes) are one-hot encoded to avoid imposing artificial numeric orders. Any plot
with a missing value in a required feature is removed entirely, across all of its survey years,
so no model ever sees an incomplete trajectory. For the terrain-conditioned RQ1 models this
removes 39 plots (156 rows) from 4survey and none from 6survey. A separate, larger loss affects
the broader candidate pool used only for feature screening, not any trained model.

\section{Environmental feature-set construction}

Three questions each need their own screened list of candidate environmental/management
variables: RQ1's predictive tiers, RQ2's shared-curve attribution set, and RQ3's plot-level
attribution set. All three follow the same recipe below, differing only in candidate pool and
target column. A small set of stand/management columns is present in every set as a protected
baseline, exempt from removal at every screening stage below (named in full where Set1 is
defined). Full per-variable diagnostics, VIF tables, and complete final column
lists are given in Appendix~\ref{app:feature-sets}; this section states the rules only.

\subsection{Deduplication, ranking and multicollinearity screening}

\begin{enumerate}
\setlength\itemsep{0pt}
\item \textbf{Coverage screening.} A few candidate variables have gaps in their source raster
  itself, not in the survey data, so some plots have no value at all for that variable. Any such
  plot is dropped the same way as described for preprocessing above: the whole plot is removed,
  which shrinks that candidate's usable rows before ranking runs. Before any feature set is
  chosen, this removes 430 plots (1,720 rows) from 4survey's full candidate pool, mostly from two
  sources: SoilGrids pH (413 plots missing) and the CEH Topographic Wetness Index (39 plots,
  mostly clustered in one compartment, consistent with a genuine gap in that raster rather than a
  coverage-boundary issue). 6survey loses none, and neither does any feature set that excludes
  these two columns.
\item \textbf{Deduplication.} Deterministic duplicates (a variable that is a documented closed-form
  function of another candidate already in the pool) are dropped outright; near-exact empirical
  duplicates (Spearman $|\rho|\geq0.95$) are reduced to one representative per pair, preferring an
  externally measured variable over a derived one.
\item \textbf{Importance ranking.} Three signals are computed against each avenue's own target and
  combined by rank-averaging: univariate Spearman $|\rho|$; permutation importance; and
  drop-column ablation (the held-out $R^2$ change from removing one variable, given every other
  candidate already in the model). Each signal has a known blind spot. Correlation ignores
  context. Permutation is biased under correlated predictors. Ablation depends on one model/split.
  This is a soft preference for cross-signal support, not a hard multi-signal gate: a variable that
  scores strongly on one signal but weakly on the others can still reach a middling combined rank.
\item \textbf{Variance-inflation screening.} Iterative VIF reduction (threshold 5.0
  \citep{dormann2013}) is applied on top of ranking, with the protected baseline columns exempt
  from removal (they remain in the design matrix, affecting every other column's VIF, but are
  never themselves dropped).
\end{enumerate}

% TODO: state explicitly which partition (training only, vs. the full pre-split pool) the
% ranking/VIF screening was computed on; whether ranking was repeated separately inside each
% spatial fold or computed once; and whether Set1-Set5 were frozen before the final evaluation
% split was drawn, so a reader can confirm the test compartments had no influence on feature-set
% construction.

\subsection{Research-question-specific feature sets}

\begin{table}[!htbp]
\centering
\scriptsize
\renewcommand{\arraystretch}{1.2}
\begin{tabular}{p{1cm}p{5.3cm}p{7.2cm}}
\toprule
Set & Selection rule & RQ1 membership (worked example) \\
\midrule
Set1 & Baseline only. & 4 columns: CanopyCover, time\_since\_thinning,
time\_since\_thinning\_missing, recent\_thinning\_5yr. \\
Set2 & Baseline + the top 10 candidates by combined importance rank (Spearman + permutation
importance + drop-column ablation, averaged); a candidate that fails VIF screening is skipped and
replaced by the next-ranked candidate, so the count stays at 10. & 14 columns: Set1 +
gwa\_weibull\_k\_50m, dist\_to\_scpt\_boundary, elevation, tas\_mean, eastness, windward\_topex,
gwa\_weibull\_a\_10m, gwa\_weibull\_a\_50m, slope\_degrees, cpmt\_compactness\_ratio. \\
Set3 & Baseline + the upper half (by combined rank) of the terrain/wind category only, ranked
within that category, not against every other candidate. VIF-screened. & 15 columns: Set1 +
gwa\_weibull\_k\_50m, elevation, eastness, windward\_topex, gwa\_weibull\_a\_10m,
gwa\_weibull\_a\_50m, slope\_degrees, gwa\_weibull\_k\_10m, solar\_radiation\_index, ceh\_twi,
gwa\_wind\_speed\_10m. \\
Set4 & Set3 + the upper half (by combined rank) of every other category: climate, soil/site,
spatial-position/edge-effects. VIF-screened. & 18 columns: Set3 + dist\_to\_scpt\_boundary,
tas\_mean, chelsa\_gdd5\_degc, cpmt\_compactness\_ratio. \\
Set5 & Baseline + every deduplicated candidate, with no importance-rank filter at all. Essentially
the full candidate pool, cleaned only of duplicates and VIF-redundant members. & 30
columns: baseline + 26 further candidates spanning every category (e.g. slope\_degrees,
soilgrids\_ph, chelsa\_gdd5\_degc, dist\_to\_road, whcl); full list in
Appendix~\ref{app:feature-sets}. \\
\bottomrule
\end{tabular}
\caption{The five feature sets created by applying the screening steps above, shown here for
RQ1's own candidate pool as a worked example. RQ2
and RQ3 apply the identical recipe over their own candidate pools (different available variables,
different target column), producing different actual membership at each Set. Full membership
for all three research questions is given in Appendix~\ref{app:feature-sets}. Baseline columns
(Set1) are protected from removal at every later stage; a fifth baseline candidate, thinning
fraction, was confirmed to be an exact linear complement of the missingness flag and is not
included separately.}
\label{tab:feature-sets-composition}
\end{table}

\begin{table}[!htbp]
\centering
\scriptsize
\renewcommand{\arraystretch}{1.2}
\begin{tabular}{p{1.3cm}p{1.3cm}p{3.6cm}p{1.4cm}p{1.6cm}p{3.4cm}}
\toprule
Research question & Set5 size & Feeds & Model comparison & Attribution & Purpose \\
\midrule
RQ1 & 30 & DNN and Chapman--Richards-guided network variants (environment-conditioned) & Primary &
No & DNN-vs-guided-network comparison under progressively richer terrain/wind/climate input; no
coefficient-based model reads these lists. \\
RQ2 & 28 & Elastic Net, XGBoost, two-stage mixed-effects model & Secondary & Primary & Which
conditions align with persistent departure from the shared curve. \\
RQ3 & 31 & Elastic Net, XGBoost, GNNWR & Secondary & Primary & Which conditions align with
plot-specific deviation; GNNWR additionally tests whether that association is spatially
varying. \\
\bottomrule
\end{tabular}
\caption{Purpose of the feature sets listed in Table~\ref{tab:feature-sets-composition}, by
research question. Each runs the identical Set1--Set5 recipe over its own candidate pool; "Set5
size" gives the total column count of each pool's largest set (baseline + every deduplicated,
VIF-screened candidate), so the differing pool sizes are visible without opening the appendix. The
model comparisons actually evaluated for all three research questions cover Set2--Set4; Set1
(baseline only) and Set5 (the full, ungated pool) are constructed by the same recipe but were not
carried through to a full model comparison. Full column lists for all three research questions:
Appendix~\ref{app:feature-sets}.}
\label{tab:feature-set-purpose}
\end{table}

% FIGURE PLACEHOLDER (pipeline overview -- simplified, high level only)
% Caption: "From plots to results: cohorts feed spatial/temporal partitions and a shared
% screening process producing Set1-5 per research question. Three branches follow: RQ1's
% predictive models, RQ2's shared-curve residual attribution, RQ3's plot-specific asymptote
% attribution, each with its own evaluation output."
%
% Confirmed evaluation output per branch (do not claim more than this without checking the
% notebooks under notebooks/spatial_analysis/, which were not inspected for this pass): RQ1
% produces per-fold and pooled predictive metrics under each split type; RQ2 produces a
% fixed-effects coefficient table plus held-out residual predictions; RQ3 produces per-fold
% predictive metrics, a per-plot coefficient table (GNNWR), and a residual spatial-autocorrelation
% diagnostic. No rendered coefficient/residual map was confirmed anywhere in the modelling
% pipeline -- do not draw one as an output box unless separately verified.
%
% Mermaid sketch (simplified, for redrawing):
% ```mermaid
% flowchart TD
%     A["Cohorts: 4survey / 6survey"] --> B["Spatial and temporal<br/>partitions"]
%     A --> C["Candidate environmental<br/>and management variables"]
%     C --> D["RQ-specific feature<br/>screening: Set1-5"]
%     B --> E1
%     B --> E2
%     B --> E3
%     D --> E1["RQ1: top-height<br/>prediction models"]
%     D --> E2["RQ2: shared-curve<br/>residual attribution"]
%     D --> E3["RQ3: plot-specific<br/>asymptote attribution"]
%     E1 --> F1["Predictive metrics,<br/>per-fold and pooled"]
%     E2 --> F2["Fixed-effects coefficient<br/>table + held-out residuals"]
%     E3 --> F3["Predictive metrics +<br/>per-plot coefficient table"]
% ```
\begin{figure}[!htbp]
\centering
\fbox{\parbox{0.9\textwidth}{\centering\vspace{4cm}FIGURE PLACEHOLDER --- pipeline overview, see
comment above source for exact content\vspace{4cm}}}
\caption{Overview of the full pipeline from raw data to the three research questions' results.}
\label{fig:pipeline-overview}
\end{figure}

\section{Top-height prediction models}

RQ1 is not a single model competition. It tests, together: predictive accuracy of the
top-height models; whether including environmental/terrain data beyond the original dataset
improves prediction (this tells stakeholders where to invest further high-quality data
acquisition); whether Chapman--Richards guidance improves prediction over an unguided network of
the same architecture; and whether any of those conclusions hold under spatial and temporal
generalisation, not only under a random/plot-level split.

Raw predictive accuracy is not the only criterion this comparison is built around. Only the
Chapman--Richards curve and its guided-network variants produce a fitted asymptote and growth rate
as an explicit, biologically interpretable output, rather than an opaque row-by-row prediction;
random forest and XGBoost have no equivalent internal structure to report, whatever their own
accuracy. This is why the guided-network family stays part of the comparison even where a
conventional baseline predicts height at least as well: they are answering a different question,
not only competing on the same one.

Target: the 95th-percentile LiDAR-derived top height, per plot per survey. Six model families
appear below. Five (the Chapman--Richards curve, linear/random-forest/XGBoost, and the DNN/PINN
family) share splits, metrics, and, where applicable, network architecture throughout.
Average-by-age is a single floor reference evaluated under one split only, not a full peer
comparison sharing every split and metric with the other five.

\subsection{Reference and conventional comparison models}

\textbf{Chapman--Richards (CR) curve.} A three-parameter curve
\begin{equation}
H(a) = y_{\max}\left(1-e^{-ka}\right)^{p}
\label{eq:cr}
\end{equation}
is fitted by nonlinear least squares (five restarts, lowest sum-of-squared-error kept), where
$H(a)$ is predicted top height at age $a$ (years), $y_{\max}$ is the asymptotic maximum height,
$k$ controls growth rate, and $p$ is a shape parameter. This is a single \emph{pooled},
population-average curve: one fit per cohort, not per plot or per compartment. It is fitted on the
training partition of the split in use (never on validation/test rows), refitted separately per
fold under the pooled spatial-block $k$-fold split so that a given fold's held-out compartments
never influence its own frozen curve. The pooled curve's fitted $(y_{\max},k,p)$ are treated as
constants downstream in the physics loss, never as trainable parameters.

\textbf{Average-by-age.} A lookup table of mean observed height per integer age year, computed on
the same training rows as the CR fit; an unseen age falls back to the overall training mean
height.

\textbf{Linear, random forest, and XGBoost baselines.} Three conventional supervised models are
fit on the same feature sets as every other RQ1 model (no-environment and, separately, each
environmental Set), evaluated under the same splits.

Linear regression fits one global coefficient per feature, assuming no interaction or curvature at
all. It sets the floor: a more complex model that cannot beat it is not learning anything beyond a
straight-line relationship between age, stand variables, and height.

Random forest averages many decision trees, each trained on a bootstrap-resampled subset of rows
and a random subset of features at each split. It can capture nonlinear interactions between
predictors with no parametric growth-curve assumption at all, giving it a genuinely different
inductive bias from the Chapman-Richards curve and the DNN/PINN family: the comparison tests
whether growth prediction needs a smooth, curve-shaped model or whether flexible, non-parametric
averaging does just as well.

XGBoost also fits an ensemble of trees, but sequentially rather than in parallel: each new tree is
trained on the residual error left by the trees already fit (gradient boosting), correcting
mistakes iteratively rather than averaging independently-grown trees as random forest does.
Gradient-boosted trees are widely used as a strong, hard-to-beat baseline in applied machine
learning precisely because this sequential correction often reaches a given accuracy with fewer,
shallower trees than random forest needs; that strength is also why it is included here, as a
demanding comparison point rather than a token baseline. The same sequential correction that makes
it strong also makes it more prone to overfitting than random forest's averaging if left
unconstrained, which is why its hyperparameters cannot simply be left at library defaults. L2
regularisation is applied at XGBoost's library default strength ($\lambda=1$); no L1 penalty is
applied ($\alpha=0$), so tree weights are shrunk towards zero rather than driven to exact
sparsity.

Its remaining hyperparameters (500 boosting rounds, maximum tree depth 4, learning rate 0.04) are
a fixed configuration, not a grid search: they reuse the setting already established for RQ3,
where they were chosen to correct a documented overfitting problem on 6survey's smaller sample,
applied here with early stopping against a held-out validation fold rather than a fixed round
count. This configuration, not XGBoost's own library defaults, is what is carried forward into the
main comparison: a direct check against those defaults, run on both the no-environment and every
environmental-feature-set variant, showed the reused configuration fits at least as well while
correcting the same overfitting risk the defaults leave unaddressed on the smaller cohort. Full
comparison: Appendix~\ref{app:hyperparameters}.

A single shared architecture is used for every DNN and Chapman--Richards-guided variant in this
project, including the environment-conditioned versions of both (Figure~\ref{fig:model-architecture}):
environmental input never grows a second trunk, only a small additional sub-network feeding the
guided variants' physics constraint. Reusing one architecture throughout means any difference
between the DNN and the guided network's results is attributable to the loss function, not to a
different network. This is a checked fair comparison, not an assumed one: a screen against four
alternative sizes was run separately on the environment-conditioned DNN and on the
Chapman--Richards-guided variants, and for neither model family did any alternative architecture
reliably improve on this same shared architecture across both cohorts at once, with the guided
variants showing even less architecture sensitivity than the DNN. Both model families
independently point to the same architecture, which is why it is retained for every model in this
chapter, not carried over untested from one family to the other.

% FIGURE PLACEHOLDER (model architecture)
% Caption: "Shared DNN/guided-network architecture. The plain data loss (DNN, and the guided
% network's own data term) uses the main trunk's output directly. The guided network adds two
% consistency terms against the Chapman-Richards curve's analytical derivative: an instantaneous
% check via automatic differentiation (physics loss) and a finite-difference check across real
% survey pairs (trajectory loss). In the environment-conditioned variants, terrain/wind features
% feed only the small asymptote (and, in the rate-conditioned variant, rate) sub-networks --
% never the main trunk -- producing a per-plot adjustment to the CR curve's asymptote/rate that
% the physics and trajectory losses are then computed against."
% Content: a block diagram. Left: [Age, no-env features] -> main trunk (3x128, shared) ->
% [predicted height]. Below/beside: [Terrain/wind features] -> small asymptote sub-network (16
% units) -> [y_max adjustment, additive] and, dashed/optional for the rate-conditioned variant,
% -> rate sub-network -> [k adjustment, multiplicative, log-space]. A separate branch shows
% [predicted height] and [age] feeding into an autograd box labelled dH/da, compared against a
% box labelled "CR analytical derivative (frozen y_max, k, p)" to produce the physics loss; a
% second small diagram shows two real survey observations of one plot feeding two forward
% passes, differenced and compared to the same CR derivative at mid-age, to produce the
% trajectory loss. These survey pairs are pre-built once from real consecutive-survey rows before
% training starts (not synthesised per batch) -- draw them as a small fixed table/lookup feeding
% the diagram, not as an on-the-fly sampling step.
\begin{figure}[!htbp]
\centering
\fbox{\parbox{0.9\textwidth}{\centering\vspace{4cm}FIGURE PLACEHOLDER --- model architecture, see
comment above source for exact content\vspace{4cm}}}
\caption{Shared DNN/Chapman--Richards-guided architecture and the two physics-consistency loss
terms.}
\label{fig:model-architecture}
\end{figure}

\subsection{Deep neural network}

Three hidden layers of 128 units each, leaky-ReLU activation, age concatenated with every other
input feature. The DNN's loss is plain mean squared error,
\begin{equation}
\mathcal{L}_{\text{data}} = \frac{1}{N}\sum_{i=1}^{N}\left(\hat H_i - H_i\right)^2,
\label{eq:data-loss}
\end{equation}
where $\hat H_i$ is the network's predicted height for row $i$ and $H_i$ is the observed height,
plus an L1 weight penalty ($\lambda_{L1}=10^{-5}$). The environment-conditioned DNN is
architecturally identical to the no-environment DNN; its terrain/wind block is simply concatenated
into the same input, widening that vector rather than adding a second network.

\subsection{Chapman--Richards-guided neural network}

Three physics-informed variants share the DNN's exact network class, differing only in loss terms
and, for the environment-conditioned variants, in an added small sub-network. The no-environment
variant uses the CR curve's $(y_{\max},k,p)$ as the single pooled, frozen cohort constants from
Equation~\ref{eq:cr}.

\paragraph*{Physics and trajectory losses.}
Every guided variant optimises the same three-term objective, using $H'(a)$, the CR curve's
analytical derivative,
\begin{equation}
H'(a) = y_{\max}\,p\,k\,e^{-ka}\left(1-e^{-ka}\right)^{p-1}.
\label{eq:cr-derivative}
\end{equation}
The \emph{physics} term matches the network's own instantaneous growth rate, obtained by automatic
differentiation of its prediction with respect to age, against the analytical rate at that same
age:
\begin{equation}
\mathcal{L}_{\text{phys}} = \frac{1}{N}\sum_{i=1}^{N}\left(\frac{\partial \hat H_i}{\partial a_i} -
H'(a_i)\right)^2.
\label{eq:physics-loss}
\end{equation}
Because age is standardised for network input, the automatic-differentiation
derivative is with respect to standardised age; it is rescaled by the training-split age standard
deviation before comparison against $H'(a_i)$, which is defined in real age units.

The \emph{trajectory} term is a separate, finite-difference check: for $M$ pairs of real,
consecutive same-plot survey observations (both endpoints drawn only from that plot's own
training-labelled rows), the network's own predicted height \emph{change} between the pair must
match the analytical rate at the pair's mid-age,
\begin{equation}
\mathcal{L}_{\text{traj}} = \frac{1}{M}\sum_{j=1}^{M}\left(\frac{\hat H_{j,\text{later}} -
\hat H_{j,\text{earlier}}}{\Delta a_j} - H'(\bar a_j)\right)^2,
\label{eq:trajectory-loss}
\end{equation}
where $\Delta a_j$ is the elapsed time between the pair's two surveys and $\bar a_j$ is their
mid-age. The two terms are deliberately distinct checks: an instantaneous derivative at one point
versus a measured change across two real observations. The full training objective is
\begin{equation}
\mathcal{L} = \mathcal{L}_{\text{data}} + \lambda_{\text{phys}}\mathcal{L}_{\text{phys}} +
\lambda_{\text{traj}}\mathcal{L}_{\text{traj}} + \lambda_{L1}\sum_{\theta}|\theta|,
\label{eq:total-loss}
\end{equation}
with $\lambda_{\text{phys}}=\lambda_{\text{traj}}=1.0$ as the default weighting (the physics-weight
ablation follows below). $(y_{\max},k,p)$ are always plain
constants at every point in this objective, never tensors that accumulate a gradient. Training
updates only the network (and, where present, the sub-networks below), never the process model
itself. A trajectory pair is only formed from two rows both labelled training under the split in
use, so no validation or test survey year can leak into this loss term.

\paragraph*{Environment-conditioned growth parameters.}

\begin{itemize}
\setlength\itemsep{0pt}
\item The terrain-conditioned variant feeds terrain/wind features into a small sub-network
  (16 hidden units) that outputs a per-plot \emph{additive} adjustment to the pooled $y_{\max}$:
  \begin{equation}
  y_{\max}^{(i)} = y_{\max} + f_{y}\!\left(\mathbf{x}^{\text{terrain}}_i\right),
  \label{eq:ymax-adjust}
  \end{equation}
  where $f_y$ is the sub-network and $\mathbf{x}^{\text{terrain}}_i$ is plot $i$'s terrain/wind
  feature vector. $k$ and $p$ remain the pooled global constants. This sub-network never sees age
  or the no-environment features, and its output is initialised near zero so training starts from
  the known-good pooled curve rather than an arbitrary value. Terrain/wind conditions only the
  growth \emph{ceiling} in this design, not the trajectory shape. This matches the
  asymptote-varies, shape-fixed framing of site-index modelling
  \citep{socha2021}. % TODO: citation -- Socha et al. 2021 (ADA/GADA), no .bib entry yet.
\item A second variant adds a further sub-network producing a per-plot \emph{multiplicative},
  log-space adjustment to $k$:
  \begin{equation}
  k^{(i)} = k \cdot \exp\!\left(f_{k}\!\left(\mathbf{x}^{\text{terrain}}_i\right)\right),
  \label{eq:k-adjust}
  \end{equation}
  again initialised near zero so $k^{(i)}\approx k$ at the start of training. $p$ stays a pooled
  global constant in every variant; no variant makes the curve's shape parameter site-specific.
\end{itemize}

\subsection{Physics-loss controls and ablation}

A physics-weight ablation sweeps $\lambda_{\text{phys}}=\lambda_{\text{traj}}\in\{0,1,2\}$ against
the default weighting, testing whether Chapman--Richards guidance itself improves prediction
rather than only whether a model labelled as guided performs well, against the matched DNN of
identical architecture and input, described above. The default ($\lambda=1.0$) is fixed at
Equation~\ref{eq:total-loss}'s own symmetric weighting before the ablation is run, not selected
afterwards from whichever value the same ablation scores best: choosing a hyperparameter from the
same test used to report it would be circular. Every guided model in the main comparison is
trained at this pre-specified default. Settings and full results: Results chapter, not repeated
here.

\begin{table}[!htbp]
\centering
\scriptsize
\renewcommand{\arraystretch}{1.2}
\begin{tabular}{p{4.6cm}p{4.4cm}p{4.4cm}}
\toprule
Model & Environmental input & Loss terms \\
\midrule
CR (baseline) & None & N/A (curve fit) \\
Average-by-age & None & N/A (lookup) \\
Linear / random forest / XGBoost & No-environment feature set & Model-specific \\
DNN (no environment) & None & Data (Eq.~\ref{eq:data-loss}) + L1 \\
DNN (environment-conditioned) & Terrain/wind, concatenated into main input & Data + L1 \\
Guided network (no environment) & None & Data + physics + trajectory + L1 \\
Guided network (terrain-conditioned asymptote) & Terrain/wind, asymptote sub-network only & Data +
physics + trajectory + L1 \\
Guided network (terrain-conditioned asymptote and rate) & Terrain/wind, asymptote and rate
sub-networks & Data + physics + trajectory + L1 \\
\bottomrule
\end{tabular}
\caption{Model family index for top-height prediction (RQ1). Full hyperparameters:
Appendix~\ref{app:hyperparameters}.}
\label{tab:rq1-models}
\end{table}

Both avenues below start from the single pooled Chapman--Richards curve fitted in
Equation~\ref{eq:cr} but diverge in what they attribute, and both report association, not causal,
relationships. RQ2 pools every plot's residual against that one shared curve, and splits into two:
RQ2a reuses RQ1's own fitted models directly, no new model is fitted; RQ2b fits new
coefficient-based models on a separate, plot-level target. RQ3 fits each plot its own asymptote
from only that plot's own repeated observations, described in its own section below.

\section{RQ2: Environmental attribution}

\paragraph*{RQ2a: does environment shrink RQ1's own departure from the curve.}

This avenue reuses RQ1's own already-fitted predictions directly, together with the fold-matched
Chapman--Richards curve: no model is fitted for RQ2a specifically. The
shared curve is the theoretical growth trajectory for Aberfoyle's average plot, so a plot's
departure from it marks whether that plot is systematically over- or under-predicted by the
population average. This avenue asks whether giving RQ1's environment-conditioned models
environmental information shrinks that departure, and whether the shrinkage concentrates
specifically on the plots the shared curve over- or under-predicts most, rather than improving
every row uniformly.

\paragraph*{Residual reduction.}
For each test row shared between the frozen, fold-matched CR baseline
and one of RQ1's three environment-conditioned model families (DNN, PINN, PINN$_k$), predictions
are joined on plot and survey year. Each prediction's residual, $r_i^{\text{CR}}$ or
$r_i^{\text{model}}$, is observed minus predicted height for that row; its \emph{magnitude},
$|r_i|$, is how far that prediction sits from the observed value, a value close to 0 meaning the
prediction is close to what was observed and a larger value meaning it is further off, regardless
of which side of the observation it falls on. A reduction is computed from these magnitudes,
\begin{equation}
\text{reduction}_i = \left|r_i^{\text{CR}}\right| - \left|r_i^{\text{model}}\right|,
\label{eq:residual-reduction}
\end{equation}
so a positive reduction means the model's prediction sits closer to the observed height than the
shared curve's does, and a negative reduction means the model is further off than the curve.

Reduction alone only says whether the error shrank, not which way either prediction was wrong.
That is carried by the residuals' \emph{sign}, not their magnitude: $r_i^{\text{CR}}$ and
$r_i^{\text{model}}$ are positive where that prediction under-predicts the observed height and
negative where it over-predicts. Reading sign alongside reduction distinguishes a model that moves
a plot's prediction toward the observed value, correcting the direction the shared curve got
wrong, from one that only shrinks the size of an error that still points the same way. This
re-reads each model's existing held-out test predictions against the identical fold's frozen CR
baseline, so RQ2a is a reinterpretation of RQ1's own results, not an additional modelling step.

\paragraph*{Quartile decomposition and coverage.}
Test rows are additionally split into four quartiles by the CR baseline's own absolute residual
within that fold (Q1: the curve's smallest errors; Q4: its largest), and the mean reduction is
reported separately per quartile rather than as one pooled figure. The split is purely rank-based,
an equal-count division of that fold's test rows by the CR baseline's own error, not a
distributional assumption; each quartile is a quarter of the population, not a handful of extreme
points, so this is a subgroup comparison, not outlier detection.

This decomposition exists because the overall mean reduction cannot distinguish two substantively
different patterns that could produce the same average: a model that shaves a similar amount of
error off every row regardless of how well the shared curve already fits it, consistent with added
model capacity alone rather than anything specifically environmental; or a model whose improvement
concentrates on the rows the shared curve fits worst, potentially at some cost on the rows it
already fits well, consistent with environmental conditioning capturing genuine, systematic
departure from a single population-average curve. Only the quartile breakdown, not the pooled
mean, can tell these apart. Coverage: all three environment-conditioned
model families, Set2--Set4, both cohorts, under the pooled spatial-block $k$-fold split described
earlier; the winning combination is additionally checked across the same reseeds used for RQ1's
own reseed-robustness check, described below, so RQ2a's headline pattern and RQ1's winner-reseed
check share one robustness result, not two separate ones.

\paragraph*{RQ2b: persistent departure from a shared growth curve.}
This avenue asks which environmental/management conditions align with a plot's persistent
position above or below the single shared curve, using only the 4survey cohort (6survey's 47
compartments were judged too few for a reliable spatial-block attribution split).

\paragraph*{Target construction.}
For each row $i$ of a plot's own survey history, a residual is computed against the pooled curve,
\begin{equation}
r_i = H_i - H(a_i),
\label{eq:cr-residual}
\end{equation}
and the target used for attribution is the plot's own mean residual across its $T$ observed
survey rows,
\begin{equation}
\overline{r}_{\text{plot}} = \frac{1}{T}\sum_{t=1}^{T} r_{i,t}.
\label{eq:mean-cr-residual}
\end{equation}
Age is excluded from every RQ2 candidate feature: it is the CR curve's only input and therefore
the sole construction variable of $\overline{r}_{\text{plot}}$, making it circular for this
target, consistent with the documented risk of age bias in age-structured growth/site-index
equations generally \citep{socha2016AgeBias}. % TODO: citation -- Socha, Coops & Ochal 2016, no
% .bib entry yet.
% TODO: verify -- confirm whether any plot contributing to Equation~\ref{eq:mean-cr-residual}'s
% average (computed over all of that plot's survey rows) falls in the test partition of the
% downstream attribution split below; the target is constructed before that split is drawn, which
% is expected rather than a leakage bug, but is not yet explicitly documented as checked.

\paragraph*{Global and mixed-effects attribution models.}
Three models are fitted on $\overline{r}_{\text{plot}}$ using RQ2's Set2--Set4, under the same
spatial-block design described earlier. Set1
(baseline only) was fitted separately as a targeted ablation, to check whether canopy cover's
attribution signal reflects a genuine environmental association or the reverse: canopy cover
itself partly reflecting growth already achieved. Set5 (the full, ungated candidate pool) was not
evaluated:

\begin{itemize}
\setlength\itemsep{0pt}
\item \textbf{Elastic Net}: continuous features standardised on the training split; unordered
  soil-class variables one-hot encoded; regularisation strength and $\ell_1$/$\ell_2$ mixing
  ratio selected by 5-fold cross-validation on the training partition only (never touching
  validation or test).
\item \textbf{XGBoost}: a separate model from RQ1's, fitted on RQ2's own target
  $\overline{r}_{\text{plot}}$ and RQ2's own candidate features rather than reused from RQ1;
  what carries over is only the hyperparameter recipe, not the fitted model itself (500 boosting
  rounds, maximum tree depth 4, learning rate 0.04, early stopping against a held-out validation
  fold; see the Reference and conventional comparison models subsection above for why this
  configuration, not library defaults, is used). This recipe was chosen for RQ1 to correct an
  overfitting problem tied to sample size, not to top-height prediction specifically, so reusing
  it for a different, plot-level target is a reasonable default rather than an unrelated
  transplant; a direct check on RQ2's own target confirmed the same configuration improves the
  fit here too, not only on RQ1's. RQ1, RQ2, and RQ3 now share one XGBoost recipe, each fitted
  separately on its own target.
\item \textbf{Two-stage mixed-effects model.} This provides an interpretable attribution baseline
  for RQ2b. Stage 1 is the already frozen, pooled CR curve above, removing the shared age--height
  relationship; Stage 2 is a \emph{linear} mixed model testing whether the resulting persistent
  plot-level departure is associated with the candidate environmental and management variables,
  \begin{equation}
  \overline{r}_{\text{plot}} = \boldsymbol{\beta}^{\top}\mathbf{x}_{\text{plot}} +
  u_{\text{cpmt}} + \varepsilon_{\text{plot}},
  \label{eq:mixed-effects}
  \end{equation}
  where $\boldsymbol{\beta}$ are fixed-effect coefficients on the standardised feature vector
  $\mathbf{x}_{\text{plot}}$, $u_{\text{cpmt}}\sim\mathcal{N}(0,\sigma_u^2)$ is a compartment-level
  random intercept, and $\varepsilon_{\text{plot}}$ is residual error. The fixed effects give the
  direction and magnitude of each association; the random intercept accounts for unmeasured
  conditions shared by plots within the same compartment, so that clustered plots are not treated
  as independent observations. This is a practical two-stage approximation to a full nonlinear
  mixed-effects refit of the CR curve itself, which was not implemented for this project. Fitted
  by restricted maximum likelihood (REML) for coefficient reporting; a separate variance-explained
  comparison between nested fixed-effects specifications uses maximum likelihood (ML) instead,
  since REML likelihoods are not directly comparable across models that differ in fixed-effects
  structure \citep{pinheiro2000}. % TODO: citation -- Pinheiro & Bates, no .bib entry yet.
  This model is fitted on the training partition only and has no held-out test evaluation step;
  its role is coefficient/variance reporting on the fitted training sample, not held-out
  prediction. Because the response is itself a residual from an already-fitted curve, this second
  stage does not propagate uncertainty from the Stage 1 curve fit; associations are interpreted as
  attribution of departure from that fitted curve, not as causal environmental effects.

  The random-effects structure is a random intercept only, no random slope on any candidate
  variable: confirmed against the fitting package's own documented default, which applies here
  since no variance-structure formula is set explicitly. Random slopes were not fitted, both
  because RQ2b poses no hypothesis about compartment-varying effects and because compartment
  sizes in this dataset range from single plots to over a thousand, too uneven to support stable
  random-slope estimation throughout.
\end{itemize}

Feature importance and SHAP values are also examined for this XGBoost fit, described more fully
under Model evaluation, uncertainty and interpretation later in this chapter. Gain-based
importance comes directly from the fit itself; SHAP is a separate step, computed by reloading the
already-fitted checkpoint afterward with no refitting involved, over the full population rather
than a held-out partition. This SHAP computation reads from the same corrected checkpoint as the
reported XGBoost fit itself, not an earlier, superseded one.

\section{RQ3: Plot-specific growth-curve deviation}

This avenue asks the same broad question at plot level: does a plot's \emph{own} fitted growth
ceiling deviate from what a static, external yield-class table would predict for it, and does
that deviation associate with environmental/management conditions. The deviation implemented here
is specifically in the fitted asymptote, with the curve's other shape parameters held fixed, for
three reasons. Each per-plot fit has as few as the cohort minimum of survey observations, too few
to reliably estimate more than one free parameter per plot; holding the shape parameters fixed
keeps the fit linear and closed-form rather than an unstable small-sample nonlinear optimisation.
The comparison target itself, the plot's asymptote against its yield-class table entry, is already
a statement about height ceiling, so this is the parameter the question is actually about, not a
compromise forced onto an unrelated one. It also matches the standard site-index-modelling
convention already used for RQ1's environment-conditioned growth parameters, described earlier in
this chapter: site quality is represented as a difference in asymptotic height potential, with
growth-curve shape held roughly constant within a species \citep{socha2021}. This is not a claim
about the full growth trajectory: a plot growing at a different rate or in a different-shaped
trajectory but reaching the same eventual ceiling is invisible to this target by construction.

\paragraph*{Target construction.}
For each plot, a per-plot asymptote $\hat y_{\max}^{\text{plot}}$ is fitted using \emph{only that
plot's own} repeated height/age observations. This is never pooled across plots, unlike
Equation~\ref{eq:cr}. Because the curve's shape parameters are held fixed per plot (taken from
its own Forestry Commission yield-class lookup, not refitted), height is linear in the one
remaining unknown, $y_{\max}$, so the fit is closed-form weighted least-squares through the
origin rather than an iterative optimisation. The target is the difference between this
plot-specific fitted asymptote and the asymptote implied by the plot's own static yield-class
table entry,
\begin{equation}
\Delta y_{\max}^{\text{plot}} = \hat y_{\max}^{\text{plot}} - y_{\max}^{\text{yldc}}.
\label{eq:local-ymax-diff}
\end{equation}
As with RQ2, this target is built entirely from height/age survey data and a static external
lookup table, with no environmental variable involved in its construction. Each per-plot fit uses
only that plot's own small number of repeated observations (as few as the cohort minimum), so the
fitted asymptote carries more sampling uncertainty for plots with fewer surveys, discussed as a
limitation below.

The plots with the most extreme $\Delta y_{\max}^{\text{plot}}$ values were checked before being
kept in the modelling population, not excluded by default. Each recurring worst-residual plot was
cross-referenced against an independent disturbance-flag record (clearfelling-like change,
measurement inconsistency, ambiguous disturbance) and a trajectory-instability cutoff (the top 1\%
of within-plot height change). None of the checked plots cleared either flag, so none were removed
on that basis; the extreme values instead reflect real, if unusually large, departures from the
yield-class expectation.

\paragraph*{Global and spatially varying attribution models.}
Elastic Net and XGBoost are fitted the same way as RQ2b's own models above, on RQ3's own
Set2--Set4 (Set1 and Set5 were not evaluated for RQ3, matching RQ2). A third
model, GNNWR (geographically and neurally weighted regression), additionally allows coefficients
to vary continuously over space rather than assuming one global linear relationship,
\begin{equation}
\hat y_i = \beta_0(s_i) + \sum_j \beta_j(s_i)\,x_{ij},
\label{eq:gnnwr}
\end{equation}
where $s_i$ is plot $i$'s spatial location and each $\beta_j(s_i)$ is a locally varying
coefficient rather than a single global one \citep{du2020}. % TODO: citation -- Du et al. 2020, no .bib entry yet.
Concretely, a global OLS fit on the same predictors is taken first and frozen; a spatial weighting
network then takes, for each plot, the vector of Euclidean distances from that plot to every plot
in a reference set (by default the training partition), and outputs one multiplicative weight per
predictor plus an intercept weight. Each local coefficient $\beta_j(s_i)$ is that plot's weight
times the corresponding frozen global coefficient, so the network learns a spatially varying
re-weighting of one shared global regression rather than an unconstrained per-plot model; one
shared network produces every coefficient's weight jointly, not a separate network per coefficient.
There is no explicit distance-decay kernel or neighbourhood/bandwidth step. Weighting is learned
implicitly from the distance vector. Held-out plots at evaluation are given their distance to
the training reference set, not a bespoke neighbour set of their own. The reference set can be
capped at a fixed row limit for memory reasons; the reported runs use the full training partition
uncapped, the same population every other model in this chapter uses. This tests whether the
attribution result itself is spatially non-stationary, a question the global Elastic Net/XGBoost
fits cannot answer. Gain-based importance and SHAP values are also examined here, described more
fully under Model evaluation, uncertainty and interpretation later in this chapter. Unlike RQ2b's
full-population computation, each of the 5 spatial-block folds' own XGBoost model is evaluated
only on its own held-out fold, and the resulting SHAP values are pooled across folds, so no plot's
SHAP explanation comes from a model that was fitted on that same plot.

% FIGURE PLACEHOLDER (RQ3 target construction and attribution models, incl. GNNWR mechanism)
% Caption: "RQ3's asymptote-deviation target, fitted from each plot's own repeated observations
% against a static yield-class expectation, feeds three attribution models: global Elastic Net
% and XGBoost, and GNNWR, whose spatial weighting network re-weights a frozen global OLS fit
% using each plot's distances to a training reference set, producing plot-level local
% coefficients and predictions evaluated under the spatial-block split. RQ2's parallel but
% simpler pipeline (pooled residual -> Elastic Net/XGBoost/two-stage mixed-effects model) is
% described in prose and equations only above, not duplicated in a figure here,
% since it has no comparable internal mechanism to unpack."
%
% Verified against the actual GNNWR implementation (third-party `gnnwr` package) before drawing:
% - input to the spatial weighting network is a vector of Euclidean distances from the focal
%   plot to every plot in the reference set (MinMax-scaled), NOT raw coordinates and NOT
%   coordinate differences;
% - no k-nearest-neighbour or distance-threshold step exists -- "neighbourhood" is implicit,
%   learned by the network from the full distance vector to the reference set;
% - the network does NOT output final coefficients directly: it outputs one multiplicative
%   weight per predictor (+ intercept), which is multiplied against a separately pre-fitted,
%   frozen global OLS coefficient to give the local coefficient beta_j(s_i);
% - one shared network with a multi-output head produces every coefficient's weight jointly, not
%   one network per coefficient;
% - held-out/test plots reuse the training reference set for their distance vector, not a
%   bespoke neighbour set built from their own location;
% - this is the standard GNNWR (Du et al. 2020) formulation; the reference set CAN be capped at a
%   fixed row count for memory reasons, but the reported runs use the full, uncapped training
%   partition (confirmed: run with --reference-set-size 0) -- do not draw a capped reference set;
% - no geographical kernel or bandwidth parameter exists in the implementation -- do not draw one;
% - no rendered spatial coefficient map was found in the modelling pipeline, only a per-plot
%   coefficient table and a Moran's I residual-autocorrelation diagnostic -- label the coefficient
%   output as a "per-plot coefficient table", not a "coefficient map", unless a map-producing
%   script is confirmed separately.
%
% Mermaid sketch (for redrawing, not final layout):
% ```mermaid
% flowchart TD
%     A["Repeated heights, own plot only"] --> B["Closed-form per-plot<br/>asymptote fit"]
%     C["Yield-class asymptote<br/>(static lookup)"] --> D["RQ3 target: y_max diff"]
%     B --> D
%     F["Environmental/management<br/>predictors (Set2-4)"] --> G["Elastic Net"]
%     F --> H["XGBoost"]
%     F --> M["Global OLS fit<br/>(frozen coefficients)"]
%     D --> G
%     D --> H
%     D --> M
%     P["Plot location"] --> Q["Distance to every plot in<br/>reference set (full training partition)"]
%     Q --> R["Spatial weighting network<br/>(one shared trunk)"]
%     R --> S["Per-predictor weights w_j(s_i)"]
%     M --> T["Local coefficient<br/>beta_j(s_i) = w_j(s_i) x beta_j"]
%     S --> T
%     T --> U["Plot-level prediction"]
%     T --> V["Per-plot coefficient table"]
%     G --> W["Spatial-block holdout evaluation"]
%     H --> W
%     U --> W
% ```
%
% Content (prose fallback if not drawn from the mermaid sketch above): left-to-right, target
% construction (closed-form per-plot asymptote fit vs. static yield-class asymptote) feeding
% three model boxes -- Elastic Net, XGBoost, and an expanded GNNWR box showing the
% distance-to-reference-set input, spatial weighting network, per-predictor weights, and their
% multiplication against the frozen global OLS coefficients -- all three converging on a
% spatial-block holdout evaluation box. A small annotation: "association, not causation."
\begin{figure}[!htbp]
\centering
\fbox{\parbox{0.9\textwidth}{\centering\vspace{4cm}FIGURE PLACEHOLDER --- RQ3 target construction
and attribution models, see comment above source for exact content\vspace{4cm}}}
\caption{Plot-specific asymptote-deviation target construction (RQ3) and its three attribution
models, including the GNNWR spatial weighting mechanism.}
\label{fig:rq3-attribution}
\end{figure}

\section{Model evaluation, uncertainty and interpretation}

This section states what is measured and why, not what was found (Results chapter).

\paragraph*{Predictive diagnostics (RQ1).}
Every RQ1 model is scored on residuals $e_i = H_i - \hat H_i$ over $N$ rows:
\begin{align*}
\text{MAE} &= \frac{1}{N}\sum_i |e_i|, &
\text{RMSE} &= \sqrt{\frac{1}{N}\sum_i e_i^2}, \\
R^2 &= 1 - \frac{\sum_i e_i^2}{\sum_i (H_i - \bar H)^2}, &
\text{Bias} &= \frac{1}{N}\sum_i e_i.
\end{align*}
$R^2$ is variance explained relative to a mean-only baseline (1 perfect, 0 matches the mean,
negative is worse than the mean on a held-out split); positive bias is systematic
under-prediction. Metrics are further broken down by five age bands to check whether an
aggregate score conceals age-specific weaknesses (the maturity gate is applied at plot level, so
individual rows below 30 can still appear; a poor $<30$ score reflects a small, non-representative
sample, not necessarily worse prediction). Pooled $R^2$ additionally carries a cluster-bootstrap
95\% confidence interval, resampling whole compartments rather than individual rows or plots. This
is RQ1 only; RQ2b and RQ3's own $R^2$ numbers do not currently carry one.

\paragraph*{Spatial autocorrelation (Moran's I).}
Tests whether nearby plots have more similar residuals than distant ones, against a null of no
spatial pattern (near zero: spatially random error; positive: errors cluster; negative: errors
alternate more than chance). Used two ways in this project, not one. On \textbf{RQ1}'s raw-height
residuals, it is a post-fit diagnostic of whether spatial structure remains unexplained (does not
feed back into feature selection). On \textbf{RQ3}'s category-removal check, the question shifts
from "is spatial structure gone" to "does removing this variable category raise residual
clustering," diagnosing whether that category was carrying real spatial signal; this check
currently runs for Set4/4survey only, not the full set/cohort matrix. It has not yet been computed
on RQ2b's own attribution residuals, or on RQ3's main GNNWR/EN/XGBoost residuals under the current
VIF-screened tiers (only superseded, pre-restructure numbers exist for the latter). A
semivariogram, fitted to test-partition residuals only, complements Moran's I with the distance
range over which spatial autocorrelation persists, rather than only whether it is present.

\paragraph*{Feature attribution (RQ2b, RQ3).}
Both use XGBoost's gain-based importance (a free byproduct of fitting) and SHAP: a separate,
deliberately computed post-hoc step decomposing one row's prediction into signed per-feature
contributions, the tree-ensemble equivalent of Elastic Net/the mixed-effects model's own
coefficients. Neither feeds back into feature selection. The role differs slightly by research
question. In \textbf{RQ2b}, SHAP is one of three converging global-attribution views, alongside
the mixed-effects model and Elastic Net. In \textbf{RQ3}, it is additionally the tool behind the
per-plot outlier check, since it is the only method here that explains an individual row's
prediction, not just average importance. SHAP values are pooled across folds into one mean per
feature in both research questions; fold-to-fold SHAP stability itself has not been computed for
either.

\paragraph*{Coefficient stability (RQ2b, RQ3).}
Elastic Net coefficients are reported as mean $\pm$ SD across the 5 spatial folds, for both RQ2b
and RQ3, checking whether a coefficient's sign and magnitude are stable across different held-out
compartments, not just significant in one fit. This has not been checked across seeds for either
research question: no seed sweep has been run for the attribution Elastic Net or XGBoost models.

\paragraph*{Robustness to stochastic and configuration choices.}
Several checks test whether a headline result depends on an incidental choice rather than the
mechanism being claimed:
\begin{itemize}
\setlength\itemsep{0pt}
\item \textbf{Training-seed reseed} (RQ1, RQ2a): the RQ1 winner is refit across 4 additional
  seeds; RQ2a reuses these same predictions rather than reseeding separately, so the two share one
  robustness result.
\item \textbf{Architecture-size sensitivity} (RQ1, all three guided-network variants): tests
  whether the fixed network size affects the headline comparison.
\item \textbf{Physics-weight ablation} (RQ1, PINN and PINN$_k$): tests whether the guided
  network's advantage, if any, depends on the physics/trajectory loss weighting rather than the
  mechanism.
\item \textbf{XGBoost hyperparameter sensitivity} (RQ1 baseline, RQ2b, RQ2a): RQ1's baseline and
  RQ2b's attribution model were both found to be running on library defaults rather than a
  deliberately chosen configuration; this check tests whether their headline results depend on
  that choice, by refitting all three under one consistent, non-default configuration. Outcome
  reported in the Results chapter, not here.
\item \textbf{Split-seed reseed} (RQ3, GNNWR, 6survey only): checks whether GNNWR's
  smaller-cohort result is stable across different fold assignments, given its wide per-fold
  variance there, discussed as a limitation later in this chapter.
\end{itemize}

\section{Methodological assumptions and limitations}

% Keep to 3-5 items. Draft points below, expand each into 2-4 sentences by hand.

\begin{itemize}
\setlength\itemsep{0pt}
\item \textbf{Spatial-block and buffer assumptions.} The 60\,m buffer is set from typical plot
  spacing, not from the measured spatial-autocorrelation range of the shared-curve residual;
  compartment-level holdout is the primary leakage control, the buffer a
  secondary one.
\item \textbf{Environmental scale mismatch.} Terrain/wind/soil variables are raster-derived and
  assigned to point-scale plots; the raster resolution does not necessarily match the scale at
  which these conditions act on a plot.
  % TODO: confirm raster resolution(s) and the assignment method (nearest-cell, bilinear, zonal
  % mean) against the actual extraction code before stating this as fact.
\item \textbf{Feature-selection leakage or split dependence.} Ranking and VIF screening, described
  earlier, determine Set1--Set5 before final evaluation; the open questions
  flagged there about which partition was used, and whether ranking was fold-specific, bear
  directly on how much confidence to place in the resulting sets.
\item \textbf{Uncertainty in plot-specific targets and their attribution models.} RQ3's per-plot
  asymptote target is fitted from as few as the cohort minimum of repeated
  observations per plot; plots with shorter survey histories contribute a noisier estimate of the
  same target as plots with longer histories, without that difference being explicitly weighted or
  flagged downstream. This compounds with 6survey's small compartment count: a multi-seed reseed
  of RQ3's GNNWR fit on 6survey shows its pooled $R^2$ sign-flipping across seeds for every feature
  set, centred near zero within the seed-to-seed spread. This is an underpowered result, not a
  reliable positive or negative one. 6survey's GNNWR fit should be read as inconclusive, not as
  evidence the spatially-varying-coefficient effect is genuinely weaker there than on 4survey.
\item \textbf{Association, not causation.} Both attribution avenues report statistical association
  between environmental/management conditions and departure from
  an expected growth curve; none of the designs support a causal claim.
\end{itemize}

\section{Reproducibility and implementation}

All split assignments, network weight initialisation, and Elastic Net/XGBoost cross-validation
folds use a fixed seed (42) unless explicitly swept as part of a stated robustness check (e.g. the
RQ1 winner-model reseed check, described above). Every model reads its feature list
from a single generated manifest, so a reported set's exact column membership can always be
checked against that source rather than a hand-copied list.

Detailed hyperparameter grids and selected values (Appendix~\ref{app:hyperparameters}), full
Set1--Set5 column lists and per-variable screening diagnostics (Appendix~\ref{app:feature-sets}),
and secondary sensitivity settings not reported in the main text (dropout-rate/architecture-width
sweeps, the narrow-gap temporal split's own results, Appendix~\ref{app:sensitivity}) are kept
out of this chapter and pointed to at the relevant section above rather than repeated here.

% TODO: verify -- confirm whether RQ2/RQ3's own model-fitting scripts read the same feature-set
% manifest as RQ1, or still require a separate wiring step, before this chapter's Set1-5
% description is presented as equally accurate for all three research questions.
